% ╔══════════════════════════════════════════════════════════════════════════════╗
% ║  Trading Algorítmico con Python y MetaTrader 5                              ║
% ║  Plantilla LaTeX para Pandoc · XeLaTeX                                      ║
% ╚══════════════════════════════════════════════════════════════════════════════╝

\documentclass[12pt, a4paper, twoside, openright]{book}

% ── Codificación ──────────────────────────────────────────────────────────────
\usepackage{fontspec}
\usepackage[spanish, es-noshorthands, es-noindentfirst]{babel}

% ── Fuentes del libro ────────────────────────────────────────────────────────
\setmainfont{Georgia}[Ligatures=TeX]
\setmonofont{Consolas}[Scale=0.86]
\setsansfont{Calibri}[Scale=0.96, Ligatures=TeX]

% ── Caracteres Unicode ────────────────────────────────────────────────────────
\usepackage{newunicodechar}
\newunicodechar{→}{\ensuremath{\rightarrow}}
\newunicodechar{←}{\ensuremath{\leftarrow}}
\newunicodechar{↑}{\ensuremath{\uparrow}}
\newunicodechar{↓}{\ensuremath{\downarrow}}
\newunicodechar{↔}{\ensuremath{\leftrightarrow}}
\newunicodechar{≈}{\ensuremath{\approx}}
\newunicodechar{≤}{\ensuremath{\leq}}
\newunicodechar{≥}{\ensuremath{\geq}}
\newunicodechar{≠}{\ensuremath{\neq}}
\newunicodechar{×}{\ensuremath{\times}}
\newunicodechar{±}{\ensuremath{\pm}}
\newunicodechar{∑}{\ensuremath{\sum}}
\newunicodechar{∞}{\ensuremath{\infty}}
\newunicodechar{α}{\ensuremath{\alpha}}
\newunicodechar{β}{\ensuremath{\beta}}
\newunicodechar{σ}{\ensuremath{\sigma}}
\newunicodechar{ε}{\ensuremath{\varepsilon}}
\newunicodechar{μ}{\ensuremath{\mu}}
\newunicodechar{ρ}{\ensuremath{\rho}}
\newunicodechar{₀}{\ensuremath{_{0}}}
\newunicodechar{₁}{\ensuremath{_{1}}}
\newunicodechar{₂}{\ensuremath{_{2}}}
\newunicodechar{₃}{\ensuremath{_{3}}}
\newunicodechar{₄}{\ensuremath{_{4}}}
\newunicodechar{²}{\ensuremath{^{2}}}
\newunicodechar{³}{\ensuremath{^{3}}}
\newunicodechar{·}{\ensuremath{\cdot}}
\newunicodechar{★}{\ensuremath{\star}}
\newunicodechar{✓}{\checkmark}
\newunicodechar{─}{\textendash}
\newunicodechar{═}{=}

% ── Página y márgenes ─────────────────────────────────────────────────────────
\usepackage[
  a4paper,
  top=2.8cm, bottom=2.8cm,
  inner=3.0cm, outer=2.2cm,
  headheight=15pt, headsep=12pt, footskip=28pt
]{geometry}

% ── Colores (prefijo bk para evitar colisión con palabras en español) ─────────
\usepackage{xcolor}
\definecolor{bkblue}  {HTML}{1a3a5c}
\definecolor{bkgold}  {HTML}{f0a500}
\definecolor{bklink}  {HTML}{1555a0}
\definecolor{bkcodebg}{HTML}{1e2b3a}
\definecolor{bkcodefg}{HTML}{e0eaf2}
\definecolor{bkcodebr}{HTML}{3a5a80}
\definecolor{bkgray}  {HTML}{707070}

% ── Imágenes ─────────────────────────────────────────────────────────────────
\usepackage{graphicx}
\usepackage{float}
\makeatletter
\def\maxwidth{\ifdim\Gin@nat@width>\linewidth\linewidth\else\Gin@nat@width\fi}
\makeatother
\setkeys{Gin}{width=\maxwidth, keepaspectratio}

% ── Hipervínculos ─────────────────────────────────────────────────────────────
\usepackage[
  unicode=true,
  colorlinks=true,
  linkcolor=bkblue,
  urlcolor=bklink,
  citecolor=bkblue,
  breaklinks=true
]{hyperref}
\hypersetup{
  pdftitle={Trading Algoritmico con Python y MetaTrader 5},
  pdfauthor={Sebastian Ospina Valencia}
}
\usepackage{url}
\urlstyle{same}

% ── Encabezados y pies — \nouppercase protege los colores ────────────────────
\usepackage{fancyhdr}
\pagestyle{fancy}
\fancyhf{}
\fancyhead[LE]{\small\color{bkgray}\textit{\nouppercase{\leftmark}}}
\fancyhead[RO]{\small\color{bkgray}\textit{\nouppercase{\rightmark}}}
\fancyfoot[LE,RO]{\small\color{bkblue}\thepage}
\renewcommand{\headrulewidth}{0.3pt}
\renewcommand{\headrule}{%
  \vskip -0.4pt\color{bkgold!60}\hrule width\headwidth height\headrulewidth\vskip 0pt}

\fancypagestyle{plain}{%
  \fancyhf{}%
  \fancyfoot[LE,RO]{\small\color{bkblue}\thepage}%
  \renewcommand{\headrulewidth}{0pt}%
}

% ── Estilo de capítulos ───────────────────────────────────────────────────────
\usepackage{titlesec}

% \protect\color{} evita que \MakeUppercase rompa el nombre del color
\titleformat{\chapter}[display]
  {\normalfont\sffamily\bfseries\protect\color{bkblue}}
  {%
    \vspace{-8pt}%
    \textcolor{bkgold}{\rule{\linewidth}{2.5pt}}\\[-4pt]%
    \hfill\fontsize{72}{72}\selectfont\color{bkgold!25}\thechapter%
  }
  {-18pt}
  {\fontsize{20}{24}\selectfont\raggedright}
  [{\vspace{4pt}\textcolor{bkgold}{\rule{\linewidth}{1pt}}}]

\titlespacing*{\chapter}{0pt}{0pt}{24pt}

\titleformat{\section}
  {\normalfont\large\bfseries\protect\color{bkblue}}
  {\color{bkgold}\thesection}
  {0.8em}{}
  [{\color{bkblue!20}\vspace{1pt}\hrule height 0.4pt}]
\titlespacing{\section}{0pt}{16pt}{6pt}

\titleformat{\subsection}
  {\normalfont\normalsize\bfseries\protect\color{bkblue}}
  {\color{bkgold!70}\thesubsection}{0.7em}{}
\titlespacing{\subsection}{0pt}{10pt}{3pt}

\titleformat{\subsubsection}
  {\normalfont\small\bfseries\itshape\protect\color{bkblue!70}}
  {}{0em}{}
\titlespacing{\subsubsection}{0pt}{8pt}{2pt}

% ── Tabla de contenidos ───────────────────────────────────────────────────────
\usepackage{tocloft}
\setlength{\cftbeforechapskip}{5pt}
\renewcommand{\cftchapfont}{\bfseries\color{bkblue}}
\renewcommand{\cftchappagefont}{\bfseries\color{bkblue}}
\renewcommand{\cftsecfont}{\small\color{bkgray}}
\renewcommand{\cftsecpagefont}{\small\color{bkgray}}
\renewcommand{\cftchapleader}{\color{bkgold!50}\cftdotfill{\cftdotsep}}
\setcounter{tocdepth}{1}

% ── Código fuente — pandoc breezedark highlighting ────────────────────────────
% Nota: fancyvrb y color los carga pandoc en 

% Caja oscura para código (mdframed)
\usepackage{mdframed}

\mdfdefinestyle{codestyle}{
  backgroundcolor  = bkcodebg,
  linecolor        = bkcodebr,
  linewidth        = 0.7pt,
  innerleftmargin  = 8pt,
  innerrightmargin = 8pt,
  innertopmargin   = 5pt,
  innerbottommargin= 5pt,
  skipabove        = 6pt,
  skipbelow        = 6pt,
  roundcorner      = 3pt
}

% Macros de resaltado de pandoc (define Shaded, NormalTok, KeywordTok, etc.)


% Redefinir Shaded con la caja oscura.
% Usa \@ifundefined por si  estaba vacío (sin bloques de código).
\makeatletter
\@ifundefined{Shaded}{%
  \newenvironment{Shaded}%
    {\begin{mdframed}[style=codestyle]\color{bkcodefg}}%
    {\end{mdframed}}%
}{%
  \renewenvironment{Shaded}%
    {\begin{mdframed}[style=codestyle]\color{bkcodefg}}%
    {\end{mdframed}}%
}
\makeatother

\newcommand{\passthrough}[1]{\texttt{\color{bkblue!80}\small#1}}

% ── Tablas ────────────────────────────────────────────────────────────────────
\usepackage{longtable}
\usepackage{booktabs}
\usepackage{array}
\usepackage{multirow}
\setlength{\tabcolsep}{8pt}
\renewcommand{\arraystretch}{1.30}

% ── Listas ────────────────────────────────────────────────────────────────────
\usepackage{enumitem}
\setlist[itemize,1]{
  label=\textcolor{bkgold}{\textbullet},
  leftmargin=1.4em, itemsep=2pt, parsep=0pt}
\setlist[itemize,2]{
  label=\textcolor{bkblue!50}{\textendash},
  leftmargin=1.4em, itemsep=1pt}
\setlist[enumerate,1]{leftmargin=1.8em, itemsep=2pt, parsep=0pt}

\providecommand{\tightlist}{%
  \setlength{\itemsep}{2pt}\setlength{\parskip}{0pt}}

% ── Citas en bloque (epígrafes) ───────────────────────────────────────────────
\mdfdefinestyle{quotestyle}{
  leftline=true, rightline=false, topline=false, bottomline=false,
  linewidth=3pt, linecolor=bkgold,
  backgroundcolor=bkblue!5,
  innerleftmargin=14pt, innerrightmargin=10pt,
  innertopmargin=8pt, innerbottommargin=8pt,
  skipabove=10pt, skipbelow=10pt
}

\renewenvironment{quote}{%
  \begin{mdframed}[style=quotestyle]%
  \itshape\color{bkblue}%
}{%
  \end{mdframed}%
}

% ── Párrafos y tipografía ─────────────────────────────────────────────────────
\usepackage{parskip}
\setlength{\parskip}{7pt}
\setlength{\parindent}{0pt}

\usepackage{setspace}
\setstretch{1.12}

\usepackage[protrusion=true, expansion=false]{microtype}

% ── Matemáticas ───────────────────────────────────────────────────────────────
\usepackage{amsmath}
\usepackage{amssymb}

% ── Figuras y tablas ─────────────────────────────────────────────────────────
\usepackage{caption}
\captionsetup{
  font={small, color=bkgray},
  labelfont={bf, color=bkblue},
  labelsep=period}
\numberwithin{figure}{chapter}
\numberwithin{table}{chapter}

% ── Viudas y huérfanas ───────────────────────────────────────────────────────
\widowpenalty=10000
\clubpenalty=10000

% ── Includes de pandoc ────────────────────────────────────────────────────────

% ══════════════════════════════════════════════════════════════════════════════
\begin{document}
% ══════════════════════════════════════════════════════════════════════════════

% ── Portada ───────────────────────────────────────────────────────────────────
\begin{titlepage}
  \thispagestyle{empty}
  \newgeometry{top=0cm, bottom=0cm, left=0cm, right=0cm}
  \noindent\includegraphics[width=\paperwidth, height=\paperheight]{portada.png}
  \restoregeometry
\end{titlepage}

% ── Página de créditos ────────────────────────────────────────────────────────
\clearpage
\thispagestyle{empty}
\mbox{}\vfill
{
  \small\color{bkgray}\setstretch{1.4}
  \noindent{\color{bkblue}\textbf{Trading Algorítmico con Python y MetaTrader 5}}\\
  \textit{Del dato al robot: estrategias cuantitativas para mercados reales}

  \vspace{10pt}
  \noindent\textbf{Autor:} Sebastián Ospina Valencia\\
  \textbf{País:} Bogotá, Colombia

  \vspace{10pt}
  \noindent{\color{bkgold}\rule{\linewidth}{0.5pt}}

  \vspace{6pt}
  \noindent\textbf{Descargo de responsabilidad:} El trading con instrumentos
  financieros implica riesgo de pérdida de capital. Los ejemplos de este libro
  son exclusivamente educativos y deben probarse en cuentas demo antes de
  usar dinero real. El autor no se
  responsabilizan por decisiones de inversión basadas en este material.

  \vspace{10pt}
  \noindent\copyright\ 2025 Sebastián Ospina Valencia. Uso educativo.
}
\clearpage

% ── Tabla de contenidos ───────────────────────────────────────────────────────
\frontmatter
\pagestyle{fancy}
\renewcommand{\contentsname}{\color{bkblue}Contenido}
\tableofcontents
\clearpage

% ── Contenido principal ───────────────────────────────────────────────────────
\mainmatter

\chapter{Prefacio}\label{prefacio}

\begin{quote}
\emph{``Lo más difícil no es ganar dinero. Lo difícil es conservarlo.''}
\end{quote}

\begin{center}\rule{0.5\linewidth}{0.5pt}\end{center}

\section{Por qué este libro existe}\label{por-quuxe9-este-libro-existe}

Cuando empecé a enseñar trading algorítmico en la Universidad del
Rosario, noté algo que se repite en cada cohorte: los estudiantes no
tienen problema con las ideas. Entienden la lógica de una media móvil,
intuyen que los precios oscilan alrededor de un valor promedio,
comprenden que los mercados tienen memoria estadística. El problema no
es conceptual.

El problema es la brecha entre la idea y el código que la ejecuta a las
3 de la mañana, sin que nadie la supervise, tomando decisiones con
dinero real.

Existe abundante literatura en inglés sobre trading cuantitativo: Ernest
Chan, Stefan Jansen, Ernie Chan, los clásicos de Edward Thorp. Pero en
español --- el idioma de más de 500 millones de personas, muchas de
ellas en economías con mercados financieros vibrantes y acceso creciente
a plataformas de trading --- la brecha es enorme. Este libro intenta
cerrarla.

No es una traducción. Es un libro construido desde cero sobre el código
real de un curso que lleva varios años ejecutándose en Bogotá, con
estudiantes reales que terminaron el programa con un robot funcionando.
Cada capítulo parte de código que fue escrito, probado, corregido y
usado en producción.

\begin{center}\rule{0.5\linewidth}{0.5pt}\end{center}

\section{A quién va dirigido}\label{a-quiuxe9n-va-dirigido}

Este libro es para ti si:

\begin{itemize}
\tightlist
\item
  \textbf{Tienes conocimientos básicos de finanzas} --- sabes qué es una
  acción, un par de divisas, un stop loss --- pero nunca has programado
  una estrategia de trading.
\item
  \textbf{Sabes algo de Python} --- o estás dispuesto a aprender lo
  necesario leyendo el Capítulo 2 antes de continuar.
\item
  \textbf{Quieres construir algo real} --- no solo entender la teoría,
  sino terminar el libro con un robot que pueda conectarse a un bróker y
  operar.
\item
  \textbf{Eres estudiante, analista, trader o curioso} que quiere
  entender cómo funcionan los sistemas cuantitativos que mueven buena
  parte del volumen de los mercados modernos.
\end{itemize}

No necesitas ser matemático ni físico. Renaissance Technologies sí los
contrata; este libro no los requiere. Necesitas disciplina, curiosidad y
disposición a escribir código aunque al principio no funcione.

\begin{center}\rule{0.5\linewidth}{0.5pt}\end{center}

\section{Qué NO es este libro}\label{quuxe9-no-es-este-libro}

Antes de continuar, tres aclaraciones importantes:

\textbf{No es un libro para hacerse rico.} El objetivo no es darte una
estrategia rentable lista para usar. Ningún libro honesto puede hacer
eso. El objetivo es darte las herramientas para que \emph{tú}
construyas, pruebas y evalúes tus propias estrategias con rigor
estadístico.

\textbf{No es un curso de Python.} Python es el medio, no el fin. El
Capítulo 2 cubre exactamente lo que necesitas saber del lenguaje para
entender el resto del libro. Para una formación completa en Python
recomendamos \emph{Python Crash Course} de Eric Matthes.

\textbf{No es un libro chartista.} No encontrarás análisis de patrones
de velas, cabeza y hombros, ni niveles de Fibonacci como herramientas
principales. El enfoque es cuantitativo: estadística, matemáticas y
código. Los indicadores técnicos aparecen, pero siempre con una
justificación estadística detrás.

\begin{center}\rule{0.5\linewidth}{0.5pt}\end{center}

\section{Advertencia de riesgo}\label{advertencia-de-riesgo}

\begin{quote}
El mercado de capitales es una de las formas más riesgosas de hacer
dinero. El trading automatizado lo es aún más: un robot mal configurado
puede ejecutar cientos de operaciones en minutos antes de que tengas
tiempo de detenerlo.

Todos los ejemplos de este libro son \textbf{educativos y explicativos}.
Ninguna estrategia presentada aquí constituye asesoría financiera ni
garantiza rentabilidad. Toda operación que realices con dinero real es
\textbf{responsabilidad exclusivamente tuya}.

Si decides operar con capital real, empieza siempre en una cuenta demo.
Entiende cada línea de código antes de ejecutarla en producción.
Gestiona tu riesgo antes de pensar en rentabilidad.
\end{quote}

\begin{center}\rule{0.5\linewidth}{0.5pt}\end{center}

\section{Cómo está organizado el
libro}\label{cuxf3mo-estuxe1-organizado-el-libro}

El libro sigue la misma arquitectura de un sistema de trading real, de
adentro hacia afuera:

\begin{verbatim}
PARTE I   — Fundamentos          (Capítulos 1–3)
PARTE II  — Operaciones Básicas  (Capítulos 4–5)
PARTE III — Arbitraje Estadístico (Capítulos 6–10)
PARTE IV  — Backtesting          (Capítulos 11–12)
PARTE V   — Productivización     (Capítulos 13–15)
\end{verbatim}

\textbf{Parte I} establece el paradigma conceptual, las herramientas de
Python necesarias y la infraestructura de conexión con MetaTrader 5.

\textbf{Parte II} construye los primeros robots operativos: el sistema
de Pivot Points y el sistema Martingala. Son simples, pero completos ---
tienen lógica de entrada, gestión de la operación y cierre.

\textbf{Parte III} introduce el arbitraje estadístico: detección de
anomalías, Bandas de Bollinger, Pairs Trading y el Robot de Regresión
Lineal. Aquí el análisis se vuelve cuantitativo en serio.

\textbf{Parte IV} cubre el backtesting: cómo evaluar una estrategia
sobre datos históricos antes de arriesgar capital real, y cómo evitar
los errores más comunes (overfitting, look-ahead bias).

\textbf{Parte V} cierra el ciclo con productivización: cómo convertir un
script en un sistema que corre de forma autónoma, con gestión de riesgo
basada en el Criterio de Kelly y capacidad de desplegarse desde GitHub.

Cada capítulo termina con el código completo comentado. Puedes seguir el
libro de principio a fin, o saltar a la parte que más te interese si ya
tienes experiencia en alguna de las áreas.

\begin{center}\rule{0.5\linewidth}{0.5pt}\end{center}

\section{Las herramientas}\label{las-herramientas}

A lo largo del libro usaremos:

\begin{longtable}[]{@{}ll@{}}
\toprule\noalign{}
Herramienta & Propósito \\
\midrule\noalign{}
\endhead
\bottomrule\noalign{}
\endlastfoot
\textbf{Python 3.9+} & Lenguaje principal \\
\textbf{MetaTrader 5} & Plataforma de bróker y fuente de datos \\
\textbf{pandas} & Manipulación de datos de precios \\
\textbf{numpy} & Cálculos numéricos \\
\textbf{scikit-learn} & Modelos de regresión (Parte III) \\
\textbf{ta} & Indicadores técnicos \\
\textbf{GitHub} & Control de versiones y despliegue del robot \\
\end{longtable}

Todas son gratuitas. Las instrucciones de instalación están en el
Capítulo 3.

\begin{center}\rule{0.5\linewidth}{0.5pt}\end{center}

\section{Una nota personal}\label{una-nota-personal}

Este curso nació de una pregunta simple que varios estudiantes me
hicieron en distintos momentos: \emph{``¿Hay algo así pero en
español?''}

La respuesta honesta era no. Este libro es el intento de cambiar eso.
Está construido sobre años de iteración en el aula, errores de
producción reales, robots que funcionaron y muchos que no, y la
convicción de que el acceso al conocimiento cuantitativo no debería
depender del idioma en que nació.

Espero que al terminar este libro tengas, al menos, un robot
funcionando. Y la claridad de saber exactamente qué está haciendo y por
qué.

\emph{Sebastián Ospina Valencia} \emph{MSc Economía, Universidad de
Antioquia} \emph{Bogotá, Colombia}

\begin{center}\rule{0.5\linewidth}{0.5pt}\end{center}

\emph{Siguiente: \href{capitulo_01_paradigma_cuantitativo.md}{Capítulo 1
--- El Paradigma Cuantitativo}}

\chapter{Capítulo 1: El Paradigma
Cuantitativo}\label{capuxedtulo-1-el-paradigma-cuantitativo}

\begin{quote}
\emph{``Los mercados no son eficientes. Son simplemente difíciles de
predecir.''} --- Jim Simons
\end{quote}

\begin{center}\rule{0.5\linewidth}{0.5pt}\end{center}

\section{1.1 Una historia sobre matemáticos y
dinero}\label{una-historia-sobre-matemuxe1ticos-y-dinero}

En 1967, un matemático llamado Edward Thorp publicó \emph{Beat the
Market}, un libro en el que demostraba que era posible valorar
sistemáticamente warrants sobre acciones usando ecuaciones
diferenciales. No era análisis chartista ni intuición: era matemática
aplicada a los precios. Thorp no solo escribió el libro --- operó la
estrategia y ganó dinero real.

Dos décadas después, Fischer Black y Myron Scholes formalizaron esa
misma idea en la ecuación que llevaría sus nombres: un modelo matemático
para valorar opciones financieras. En 1997 ganaron el Premio Nobel de
Economía. Los mercados financieros habían entrado definitivamente en la
era cuantitativa.

Pero fue Jim Simons quien llevó esto al extremo. Simons era matemático
del MIT, especialista en geometría diferencial, y en 1988 fundó
Renaissance Technologies con un objetivo radical: operar los mercados
\textbf{exclusivamente} con modelos matemáticos y estadísticos, sin
ninguna intervención de analistas financieros tradicionales. Su fondo
Medallion generó retornos anualizados superiores al 60\% durante más de
tres décadas --- el mejor desempeño sostenido en la historia de los
mercados financieros.

La historia de Simons y sus colegas está narrada magistralmente en
\emph{The Man Who Solved the Market} de Gregory Zuckerman, y la historia
más amplia del fenómeno quant está en \emph{The Quants} de Scott
Patterson. Ambos libros aparecen en la bibliografía de este libro no
como referencia técnica, sino como contexto: para entender de dónde
vienen las ideas que vamos a implementar en código.

Lo que Simons demostró --- y lo que este libro busca transmitir --- es
que los mercados tienen estructura estadística explotable. No siempre.
No fácilmente. No sin riesgo. Pero sí de forma sistemática, con las
herramientas adecuadas.

Hoy esas herramientas están disponibles para cualquier persona con una
laptop, conexión a internet y disposición para aprender. Eso es, en
esencia, lo que hace posible este libro.

\begin{center}\rule{0.5\linewidth}{0.5pt}\end{center}

\section{1.2 ¿Qué es el Trading
Algorítmico?}\label{quuxe9-es-el-trading-algoruxedtmico}

\begin{quote}
\emph{Son instrumentos financieros con base en un algoritmo que permite
realizar operaciones financieras de manera automática. Su objetivo
principal es generar Alpha.}
\end{quote}

Esta definición, directa y precisa, captura los dos elementos
esenciales:

\begin{enumerate}
\def\labelenumi{\arabic{enumi}.}
\tightlist
\item
  \textbf{Un algoritmo} --- un conjunto de reglas explícitas,
  reproducibles y basadas en datos que determinan cuándo comprar, cuándo
  vender y cuánto arriesgar.
\item
  \textbf{Automatización} --- las decisiones se ejecutan sin
  intervención humana en el momento de la operación.
\end{enumerate}

La palabra clave es \emph{Alpha}.

\subsection{1.2.1 Alpha vs.~Buy and Hold}\label{alpha-vs.-buy-and-hold}

En finanzas, \textbf{Alpha} es el rendimiento que una estrategia genera
de forma \emph{independiente del mercado}. Para entenderlo, compáralo
con su opuesto:

\textbf{Buy and Hold} (\emph{comprar y mantener}) es la estrategia más
simple: compras un activo --- digamos el S\&P 500 o Bitcoin --- y lo
mantienes durante años, capturando el crecimiento general del mercado.
Si el mercado sube 10\%, tú subes 10\%. Si cae 40\%, tú caes 40\%. No
hay habilidad activa involucrada; es exposición pura al \textbf{Beta}
del mercado.

\textbf{Alpha} es el rendimiento que va \emph{más allá} del mercado. Si
el mercado sube 10\% y tu estrategia sube 18\%, el Alpha es ese 8\%
adicional. Más importante aún: si el mercado cae 30\% y tu estrategia
cae solo 5\%, o incluso genera retornos positivos --- eso es Alpha
genuino. Es rendimiento independiente de la dirección del mercado.

\emph{{[} Diagrama --- disponible en la versión digital {]}}

\textbf{Figura 1.1 --- Alpha vs Beta.} El trading algorítmico busca
retornos independientes de la dirección del mercado, no simplemente
seguirlo.

\begin{quote}
\textbf{Nota importante:} Generar Alpha consistente es
extraordinariamente difícil. La mayoría de los hedge funds no logran
superar consistentemente a sus benchmarks de largo plazo. Este libro no
promete una fórmula para lograrlo --- promete las herramientas para
\emph{buscarlo} con rigor y honestidad estadística.
\end{quote}

\subsection{1.2.2 Automatización no significa ausencia
humana}\label{automatizaciuxf3n-no-significa-ausencia-humana}

Un malentendido frecuente: que un robot de trading opera ``solo'' y no
requiere atención. La realidad es diferente:

\begin{itemize}
\tightlist
\item
  El robot ejecuta las \emph{decisiones operativas} de forma autónoma
  (comprar, vender, ajustar stops).
\item
  El ser humano se encarga del \emph{desarrollo, monitoreo y
  mantenimiento} del sistema.
\end{itemize}

Un robot mal diseñado puede ejecutar cientos de operaciones erróneas en
minutos. Un robot bien diseñado igual requiere supervisión periódica:
los mercados cambian, las condiciones estadísticas evolucionan, y ningún
modelo es válido para siempre.

\textbf{El trading algorítmico requiere una inversión seria en
desarrollo y monitoreo continuo.} Quien lo entiende como una caja negra
que genera dinero mientras duermes está en el camino equivocado.

\begin{center}\rule{0.5\linewidth}{0.5pt}\end{center}

\section{1.3 Ventajas y Desventajas}\label{ventajas-y-desventajas}

\subsection{1.3.1 Ventajas}\label{ventajas}

\begin{longtable}[]{@{}
  >{\raggedright\arraybackslash}p{(\columnwidth - 2\tabcolsep) * \real{0.5000}}
  >{\raggedright\arraybackslash}p{(\columnwidth - 2\tabcolsep) * \real{0.5000}}@{}}
\toprule\noalign{}
\begin{minipage}[b]{\linewidth}\raggedright
Ventaja
\end{minipage} & \begin{minipage}[b]{\linewidth}\raggedright
Descripción
\end{minipage} \\
\midrule\noalign{}
\endhead
\bottomrule\noalign{}
\endlastfoot
\textbf{Reducción de errores humanos} & El robot no se equivoca al
calcular un precio, no olvida colocar el stop loss, no ingresa un
volumen incorrecto por error de tipeo. \\
\textbf{Menor desgaste emocional} & El miedo y la codicia son los peores
consejeros en trading. Un algoritmo no tiene emociones: ejecuta las
reglas sin importar si el mercado está cayendo en pánico o eufórico al
alza. \\
\textbf{Señales claras y precisas} & Una estrategia algorítmica define
exactamente cuándo entrar: no ``cuando se vea fuerte'', sino ``cuando el
precio supere la media de 20 períodos por más de 1.5 desviaciones
estándar''. \\
\textbf{Metodología reproducible} & El proceso de diseño, prueba y
evaluación de estrategias cuantitativas es sistemático. Puedes
repetirlo, mejorarlo y comunicarlo a otros. \\
\textbf{Escalabilidad} & El mismo robot puede operar 1 símbolo o 50 con
el mismo esfuerzo. Un trader humano tiene límites físicos de atención;
el código no. \\
\end{longtable}

\subsection{1.3.2 Desventajas}\label{desventajas}

\begin{longtable}[]{@{}
  >{\raggedright\arraybackslash}p{(\columnwidth - 2\tabcolsep) * \real{0.5000}}
  >{\raggedright\arraybackslash}p{(\columnwidth - 2\tabcolsep) * \real{0.5000}}@{}}
\toprule\noalign{}
\begin{minipage}[b]{\linewidth}\raggedright
Desventaja
\end{minipage} & \begin{minipage}[b]{\linewidth}\raggedright
Descripción
\end{minipage} \\
\midrule\noalign{}
\endhead
\bottomrule\noalign{}
\endlastfoot
\textbf{Alto costo de desarrollo en tiempo} & Diseñar, programar, probar
y depurar un robot robusto requiere semanas o meses de trabajo. No hay
atajos. \\
\textbf{Monitoreo continuo} & Un robot que corre 24/7 necesita
supervisión: conexiones caídas, errores de ejecución, cambios en las
condiciones del bróker. \\
\textbf{Riesgo de producción} & Una mala salida a producción puede
generar pérdidas grandes y rápidas antes de que puedas intervenir. La
Parte V de este libro dedica un capítulo entero a mitigar este
riesgo. \\
\textbf{Costos de despliegue} & Servidores, conexiones VPS, cuentas de
bróker, licencias de datos --- todo tiene un costo que hay que
considerar en la ecuación de rentabilidad. \\
\textbf{Habilidades especializadas} & Requiere estadística, programación
y conocimiento de mercados. La intersección de estas tres áreas no es
trivial de alcanzar. \\
\end{longtable}

\begin{center}\rule{0.5\linewidth}{0.5pt}\end{center}

\section{1.4 Los Tres Tipos de
Estrategias}\label{los-tres-tipos-de-estrategias}

Desde el punto de vista cuantitativo, cualquier estrategia de trading
puede clasificarse en una de tres familias. Esta taxonomía es
fundamental porque determina la lógica estadística del sistema, las
condiciones de mercado en que funciona y los indicadores apropiados para
construirla.

\emph{{[} Diagrama --- disponible en la versión digital {]}}

\textbf{Figura 1.2 --- Taxonomía de estrategias cuantitativas.} Las tres
familias no son mutuamente excluyentes; muchas estrategias reales
combinan elementos de más de una.

\subsection{1.4.1 Estrategias de Momentum (o Trend
Following)}\label{estrategias-de-momentum-o-trend-following}

\textbf{Premisa:} Lo que está sucediendo en el mercado \emph{continuará
sucediendo} en el corto plazo.

Si el precio ha estado subiendo, la estrategia asume que seguirá
subiendo y compra. Si ha estado bajando, vende. Es la apuesta por la
persistencia de la tendencia.

\begin{itemize}
\tightlist
\item
  \textbf{Momento Absoluto:} ¿El precio de hoy es mayor que el de hace N
  períodos? Compra si sí, vende si no.
\item
  \textbf{Price Breakthrough:} El precio rompe un nivel de resistencia
  significativo --- señal de que la tendencia se está acelerando.
\item
  \textbf{Rate of Change (ROC):} Mide la velocidad del cambio de precio.
  Un ROC alto indica momentum fuerte.
\end{itemize}

\begin{quote}
\textbf{En qué condiciones funciona:} Mercados que están en tendencia
clara (trending markets). Falla en mercados laterales o que oscilan sin
dirección definida.
\end{quote}

\subsection{1.4.2 Estrategias de Mean Reversion (Reversión a la
Media)}\label{estrategias-de-mean-reversion-reversiuxf3n-a-la-media}

\textbf{Premisa:} Los precios oscilan alrededor de un valor promedio y,
cuando se alejan demasiado, tienden a volver.

Esta familia tiene dos preguntas fundacionales que determinan la calidad
de la estrategia: - \emph{¿A cuál media?} --- una media móvil simple,
exponencial, el precio fundamental, el valor justo estadístico. -
\emph{¿Cuándo la diferencia es estadísticamente significativa?} --- no
toda desviación del promedio es una oportunidad; hay que distinguir
ruido de señal.

\begin{itemize}
\tightlist
\item
  \textbf{Martingalas:} Duplican el tamaño de la posición después de
  cada pérdida, apostando a que el precio revertirá. Son matemáticamente
  peligrosas sin gestión de riesgo adecuada (Capítulo 5).
\item
  \textbf{Spreads deviation:} El spread entre dos activos
  correlacionados se aleja de su promedio histórico --- señal de
  reversión en el Pairs Trading (Capítulo 9).
\item
  \textbf{Detección de Anomalías:} Cuando el precio se desvía de su
  modelo estadístico de referencia por encima de un umbral --- la base
  del Robot de Regresión (Capítulo 10).
\item
  \textbf{Bandas de Bollinger:} Precio por debajo de la banda inferior →
  compra; precio por encima de la banda superior → venta (Capítulo 8).
\end{itemize}

\begin{quote}
\textbf{En qué condiciones funciona:} Mercados laterales o con reversión
estadística comprobada. Falla en mercados con tendencias prolongadas.
\end{quote}

\subsection{1.4.3 Estrategias
Predictivas}\label{estrategias-predictivas}

\textbf{Premisa:} Un modelo estadístico o matemático puede estimar los
niveles futuros del precio con suficiente precisión para generar señales
rentables.

\begin{itemize}
\tightlist
\item
  \textbf{Regresión Lineal:} Ajustar una línea al precio histórico
  reciente y operar las desviaciones del precio actual respecto a ese
  modelo (el Robot REG del Capítulo 10).
\item
  \textbf{Machine Learning:} Modelos como Random Forest o Gradient
  Boosting entrenados sobre features de mercado para predecir la
  dirección del siguiente período.
\item
  \textbf{Modelos ARIMA:} Modelos estadísticos de series de tiempo que
  capturan autocorrelación en los retornos.
\end{itemize}

\begin{quote}
\textbf{En qué condiciones funciona:} Requieren periodos de
reentrenamiento frecuente. Son susceptibles al overfitting --- el
enemigo número uno de las estrategias predictivas (Capítulo 12).
\end{quote}

\begin{center}\rule{0.5\linewidth}{0.5pt}\end{center}

\section{1.5 La Arquitectura Básica de un Sistema de
Trading}\label{la-arquitectura-buxe1sica-de-un-sistema-de-trading}

Antes de escribir una línea de código, es útil entender las siete capas
que componen cualquier sistema de trading algorítmico. Este libro las
recorre en orden:

\emph{{[} Diagrama --- disponible en la versión digital {]}}

\textbf{Figura 1.3 --- Las siete capas de un sistema de trading
algorítmico.}

\begin{longtable}[]{@{}
  >{\raggedright\arraybackslash}p{(\columnwidth - 4\tabcolsep) * \real{0.3333}}
  >{\raggedright\arraybackslash}p{(\columnwidth - 4\tabcolsep) * \real{0.3333}}
  >{\raggedright\arraybackslash}p{(\columnwidth - 4\tabcolsep) * \real{0.3333}}@{}}
\toprule\noalign{}
\begin{minipage}[b]{\linewidth}\raggedright
Capa
\end{minipage} & \begin{minipage}[b]{\linewidth}\raggedright
Descripción
\end{minipage} & \begin{minipage}[b]{\linewidth}\raggedright
Capítulos de este libro
\end{minipage} \\
\midrule\noalign{}
\endhead
\bottomrule\noalign{}
\endlastfoot
\textbf{Infraestructura} & Hardware, software, plataforma del bróker &
Cap. 3 \\
\textbf{Conexión a datos} & Autenticación con MT5, gestión de
credenciales & Cap. 3 \\
\textbf{Obtención de datos} & Extracción de velas OHLCV, ticks,
posiciones abiertas & Cap. 3 \\
\textbf{Código de estrategia} & La lógica: calcular señales, determinar
entradas y salidas & Caps. 4--10 \\
\textbf{Backtesting} & Probar la estrategia sobre datos históricos &
Caps. 11--12 \\
\textbf{Código de trading} & Enviar órdenes, gestionar posiciones,
calcular riesgo & Caps. 4--10 \\
\textbf{Productivización} & Automatizar la ejecución, logs, monitoreo,
despliegue & Caps. 13--15 \\
\end{longtable}

Un error frecuente de los principiantes es saltar directamente a la capa
4 (la estrategia) sin haber construido las capas anteriores. El
resultado es código que ``funciona en mi máquina'' pero falla en
producción porque no maneja errores de conexión, no tiene logs, o no
puede recuperarse de una caída del servidor.

Este libro respeta ese orden.

\begin{center}\rule{0.5\linewidth}{0.5pt}\end{center}

\section{1.6 El Arbitraje Estadístico: la columna vertebral del
libro}\label{el-arbitraje-estaduxedstico-la-columna-vertebral-del-libro}

La Parte III de este libro está dedicada al \textbf{Arbitraje
Estadístico} --- la familia de estrategias que fundamenta la mayoría del
código del repositorio.

\begin{quote}
\emph{``Explotar las características estadísticas del activo o los
activos con el fin de generar rentabilidades.''}
\end{quote}

A diferencia del arbitraje clásico (que busca diferencias de precio del
mismo activo en dos mercados distintos), el arbitraje estadístico no
requiere precios idénticos. Requiere que haya una \textbf{relación
estadística estable} entre variables de mercado --- una relación que el
modelo puede identificar, cuantificar y explotar cuando se desvía
temporalmente.

Los ejemplos concretos que construiremos:

\begin{itemize}
\tightlist
\item
  \textbf{Detección de Anomalías:} Cuando el precio se desvía
  significativamente de su modelo de regresión lineal reciente.
\item
  \textbf{Bandas de Bollinger:} Cuando el precio cruza los límites
  estadísticos de desviación estándar.
\item
  \textbf{Pairs Trading:} Cuando el spread entre dos activos
  cointegrados se aleja de su media histórica.
\item
  \textbf{Robot REG:} Un modelo de regresión que predice el precio
  ``justo'' y opera las desviaciones.
\end{itemize}

En todos los casos, la pregunta central es la misma: \emph{¿cuándo la
diferencia entre el precio actual y su valor estadístico de referencia
es lo suficientemente grande y persistente como para justificar una
operación?}

\begin{center}\rule{0.5\linewidth}{0.5pt}\end{center}

\section{1.7 Resumen del Capítulo}\label{resumen-del-capuxedtulo}

\begin{longtable}[]{@{}
  >{\raggedright\arraybackslash}p{(\columnwidth - 2\tabcolsep) * \real{0.5000}}
  >{\raggedright\arraybackslash}p{(\columnwidth - 2\tabcolsep) * \real{0.5000}}@{}}
\toprule\noalign{}
\begin{minipage}[b]{\linewidth}\raggedright
Concepto
\end{minipage} & \begin{minipage}[b]{\linewidth}\raggedright
Definición
\end{minipage} \\
\midrule\noalign{}
\endhead
\bottomrule\noalign{}
\endlastfoot
\textbf{Trading Algorítmico} & Sistema de reglas explícitas basadas en
datos que ejecuta operaciones financieras automáticamente \\
\textbf{Alpha} & Rendimiento generado independientemente del movimiento
general del mercado \\
\textbf{Beta} & Exposición al movimiento general del mercado (Buy and
Hold) \\
\textbf{Momentum} & Estrategia que apuesta por la persistencia de la
tendencia actual \\
\textbf{Mean Reversion} & Estrategia que apuesta por el retorno del
precio a su valor promedio \\
\textbf{Predictiva} & Estrategia que usa modelos estadísticos para
estimar niveles futuros \\
\textbf{Arbitraje Estadístico} & Explotar relaciones estadísticas
estables entre variables de mercado \\
\textbf{Overfitting} & Ajustar demasiado el modelo a datos históricos,
perdiendo capacidad predictiva real \\
\end{longtable}

\begin{center}\rule{0.5\linewidth}{0.5pt}\end{center}

\section{1.8 Lecturas Recomendadas}\label{lecturas-recomendadas}

Para profundizar en el contexto histórico y conceptual del trading
cuantitativo:

\begin{itemize}
\tightlist
\item
  \textbf{Beat the Market} --- Edward O. Thorp \& Sheen T. Kassouf
  (1967). El libro que dio inicio al trading sistemático moderno.
\item
  \textbf{The Quants} --- Scott Patterson (2010). La historia de cómo
  los matemáticos conquistaron Wall Street.
\item
  \textbf{The Man Who Solved the Market} --- Gregory Zuckerman (2019).
  La biografía no autorizada de Jim Simons y Renaissance Technologies.
\item
  \textbf{Algorithmic Trading: Winning Strategies and Their Rationale}
  --- Ernest P. Chan (2013). La referencia técnica más práctica
  disponible en inglés.
\item
  \textbf{Machine Learning for Algorithmic Trading} --- Stefan Jansen
  (2020). El estado del arte en estrategias basadas en machine learning.
\end{itemize}

\begin{center}\rule{0.5\linewidth}{0.5pt}\end{center}

\emph{En el siguiente capítulo introduciremos Python --- el lenguaje que
usaremos para implementar todo lo que se describió aquí.}

\emph{Siguiente: \href{capitulo_02_introduccion_python.md}{Capítulo 2
--- Introducción a Python para Trading Algorítmico}}

\end{document}
